\chapter{The abc Conjecture}\label{ch:abc}

\begin{center}\itshape
If $a+b=c$ with no common factor, then $c$ is rarely much larger than the product of the distinct primes dividing $abc$.\\
\upshape(Masser and Oesterl\'e, 1985, in our words)
\end{center}

Take two numbers made of small primes, say $1$ and $8=2^3$, and add them: $9=3^2$ is again made of small primes. This is rare. The sum of two
numbers usually picks up a large new prime factor, as $16+81=97$ does. Additive structure (the sum) and multiplicative structure (the factorisations)
seem to be independent, and the \emph{abc conjecture} is the most precise statement of this independence that anybody has dared to make.

It came out of the theory of elliptic curves. In 1985 Joseph Oesterl\'e, thinking about a conjecture of Lucien Szpiro on the discriminants of elliptic
curves, and David Masser, thinking about an analogy with polynomials, arrived at the same statement. It has a striking list of consequences: Fermat's
last theorem for all large exponents, Roth's theorem on rational approximation, Faltings's theorem on rational points of curves, and many more, each
following in a few lines. Dorian Goldfeld called it the most important unsolved problem in Diophantine analysis. For polynomials it is a theorem with a
one-page proof, which we give; for integers it is open. In 2012 Shinichi Mochizuki announced a proof, which most experts do not accept; we tell that story
at the end.

This chapter states the conjecture, shows the evidence, proves the polynomial version and its corollary Fermat's last theorem for polynomials, derives
consequences of abc for Fermat-type equations, and then turns to a cousin among the problems about sums and fractions, the Erd\H os--Straus conjecture
on $\frac4n$, of which we prove everything that polynomial identities can give.

\begin{think}
Find coprime $a,b$ with $a+b=c$ such that the product of the distinct primes dividing $abc$ is smaller than $c$. (One example is in the first paragraph.)
Can you find one with $c=2^k$ or $c=3^k$? How about infinitely many?
\end{think}

\section{The conjecture and the evidence}

The \emph{radical} of a positive integer $n$ is the product of its distinct prime factors, $\operatorname{rad}(n)=\prod_{p\mid n}p$; so $\operatorname{rad}(72)=6$
and $\operatorname{rad}(1)=1$. An \emph{abc triple} is a triple of positive integers with $a+b=c$ and $\gcd(a,b)=1$ (hence $a$, $b$, $c$ are pairwise
coprime, and $\operatorname{rad}(abc)=\operatorname{rad}(a)\operatorname{rad}(b)\operatorname{rad}(c)$). Its \emph{quality} is
\[
q(a,b,c)=\frac{\log c}{\log\operatorname{rad}(abc)}.
\]
A triple with $q>1$, that is $c>\operatorname{rad}(abc)$, is an \emph{abc hit}.

\begin{center}
\emph{The abc conjecture.} For every $\varepsilon>0$ there are only finitely many abc triples with $q(a,b,c)>1+\varepsilon$. Equivalently, for every
$\varepsilon>0$ there is a constant $K_\varepsilon$ with $c<K_\varepsilon\operatorname{rad}(abc)^{1+\varepsilon}$ for every abc triple.
\end{center}

(The two forms are equivalent: finitely many exceptions can be absorbed into the constant, and conversely $c<K_\varepsilon\operatorname{rad}(abc)^{1+\varepsilon}$
with $q>1+2\varepsilon$ gives $\operatorname{rad}(abc)^{1+2\varepsilon}<c<K_\varepsilon\operatorname{rad}(abc)^{1+\varepsilon}$, which bounds
$\operatorname{rad}(abc)$ and then $c$.)

The $\varepsilon$ cannot be removed.

\begin{proposition}\label{prop:abcinf}
For $k\ge1$, the triple $\bigl(1,\,3^{2^k}-1,\,3^{2^k}\bigr)$ is an abc triple with $\dfrac c{\operatorname{rad}(abc)}>\dfrac{2^{k+1}}3$. In particular there are
infinitely many abc hits, and $c/\operatorname{rad}(abc)$ is unbounded.
\end{proposition}
\begin{proof}
Let $b_k=3^{2^k}-1$. We show $2^{k+2}\mid b_k$ by induction. For $k=1$, $b_1=8$. And $b_{k+1}=\bigl(3^{2^k}-1\bigr)\bigl(3^{2^k}+1\bigr)=b_k\bigl(3^{2^k}+1\bigr)$ with
$3^{2^k}+1$ even. So $\operatorname{rad}(b_k)\le2\cdot\frac{b_k}{2^{k+2}}$, since dividing out $2^{k+2}$ and putting back one factor $2$ leaves a multiple of
$\operatorname{rad}(b_k)$. Hence $\operatorname{rad}(1\cdot b_k\cdot3^{2^k})\le3\cdot\frac{b_k}{2^{k+1}}<\frac{3\cdot3^{2^k}}{2^{k+1}}$.
\end{proof}

How large can the quality be? The record, found by \'Eric Reyssat, is
\[
2+3^{10}\cdot109=23^5,\qquad q=\frac{\log 6\,436\,343}{\log(2\cdot3\cdot109\cdot23)}=\frac{\log6\,436\,343}{\log15\,042}=1.62991\ldots,
\]
followed by Benne de Weger's $11^2+3^2\cdot5^6\cdot7^3=2^{21}\cdot23$, with $q=1.62599\ldots$ Below $10^7$ there are exactly $3504$ abc hits
(Figure~\ref{fig:abcquality}, computed for this book), and only $17$ of them have $q>1.4$. Large searches, among them the distributed project
\emph{ABC@Home}, have listed the hits much further, and Reyssat's triple is still the record as far as we know.

\begin{figure}[ht]
\centering
\fig{ch09-quality}
\caption{All $3504$ abc hits with $c<10^7$: the quality $q$ against $c$. The hits of quality above $1.4$ are in red. Hits become more frequent as $c$
grows, but high quality stays rare: the abc conjecture says that above any line $q=1+\varepsilon$ there are only finitely many points.}
\label{fig:abcquality}
\end{figure}

\paragraph{The heuristic.} Why should the conjecture be true? The integers up to $X$ whose radical is at most $X^\theta$ are few: about $X^{\theta+o(1)}$
of them (there are at most $X^\theta$ possible radicals, and each radical $r$ is the radical of few numbers up to $X$, a bound we do not prove here).
Choose $a$ and $b$ up to $X$ with radicals at most $X^\alpha$ and $X^\beta$: about $X^{\alpha+\beta}$ pairs. If additive and multiplicative structure
are independent, the chance that $c=a+b$ has radical at most $X^\gamma$ is about $X^{\gamma-1}$. So we expect about $X^{\alpha+\beta+\gamma-1}$ triples.
A triple of quality at least $1+\varepsilon$ has $\alpha+\beta+\gamma\le\frac1{1+\varepsilon}$, so the expected number at the scale $X$ is at most
$X^{-\varepsilon/(1+\varepsilon)+o(1)}$, and summing over the scales $X=2,4,8,\dots$ gives a finite total. At quality exactly $1$ the exponent is $0$,
and the infinitely many hits of Proposition~\ref{prop:abcinf} live on that edge. This is a heuristic, not a proof: the independence of addition and
multiplication is exactly what nobody can prove.

\section{Polynomials}

For polynomials over $\mathbb C$, the radical becomes the number of distinct roots. For a nonzero polynomial $f$, let $n_0(f)$ be the number of distinct
complex roots of $f$.

\begin{theorem}[Stothers, 1981; Mason, 1984]\label{thm:mason}
Let $a,b,c\in\mathbb C[t]$ be coprime polynomials with $a+b=c$, not all constant. Then
\[
\max(\deg a,\deg b,\deg c)\le n_0(abc)-1.
\]
\end{theorem}
\begin{proof}
(Snyder's proof.) Since $a$ and $b$ are coprime and $c=a+b$, the three polynomials are pairwise coprime, so they have no common roots. Consider the
\emph{Wronskian} $W=ab'-a'b$. Substituting $b=c-a$ and $a=c-b$,
\[
W=ac'-a'c=-(bc'-b'c).
\]
\emph{$W\ne0$.} Otherwise $(b/a)'=W/a^2=0$, so $b=\lambda a$ for a constant $\lambda$; as $a$ and $b$ are coprime, both are constant, and so is $c$.

\emph{Many roots of $W$.} Let $\alpha$ be a root of $a$ of multiplicity $m$. Then $(t-\alpha)^{m-1}$ divides $a$ and $a'$, hence $W=ab'-a'b$. The same holds for
the roots of $b$ (same formula) and of $c$ (formula $W=ac'-a'c$). These roots are distinct, so
\[
D=\prod_{\alpha:\ (abc)(\alpha)=0}(t-\alpha)^{m_\alpha-1}
\]
divides $W$, where $m_\alpha$ is the multiplicity of $\alpha$ as a root of $abc$. Its degree is $\deg a+\deg b+\deg c-n_0(abc)$.

\emph{Few roots of $W$.} $\deg W\le\deg a+\deg b-1$: the terms of degree $\deg a+\deg b$ in $ab'$ and $a'b$ vanish (their leading coefficients are
$\deg b\cdot\ell$ and $\deg a\cdot\ell$ times $t^{\deg a+\deg b-1}$ for the product $\ell$ of the leading coefficients, so even degree $\deg a+\deg b-1$ is
an upper bound). Hence
\[
\deg a+\deg b+\deg c-n_0(abc)\le\deg W\le\deg a+\deg b-1,
\]
that is $\deg c\le n_0(abc)-1$. Using $W=ac'-a'c$ (so $\deg W\le\deg a+\deg c-1$) gives $\deg b\le n_0(abc)-1$, and $W=-(bc'-b'c)$ gives
$\deg a\le n_0(abc)-1$.
\end{proof}

\begin{corollary}[Fermat's last theorem for polynomials]\label{cor:fermatpoly}
For $n\ge3$ there are no coprime polynomials $x,y,z\in\mathbb C[t]$, not all constant, with $x^n+y^n=z^n$.
\end{corollary}
\begin{proof}
Apply Theorem~\ref{thm:mason} to $a=x^n$, $b=y^n$, $c=z^n$, which are coprime. The roots of $x^ny^nz^n$ are the roots of $x$, $y$ and $z$, so
$n_0(x^ny^nz^n)\le\deg x+\deg y+\deg z=:d$. Then $n\deg x\le d-1$, and likewise for $y$ and $z$; adding, $nd\le3d-3<3d$, so $n<3$.
\end{proof}

The theorem is sharp (Exercise~\ref{ex:masonsharp}), and it needs characteristic $0$ only to get $W\ne0$ (Exercise~\ref{ex:charp}). It is the model for the
integer conjecture: $\deg$ plays the role of $\log c$, and $n_0$ the role of $\log\operatorname{rad}(abc)$. For polynomials the derivative does the work; the
integers have no derivative, and that is one way to say why abc is hard.

\section{What abc would give}

\begin{proposition}\label{prop:abcfermat}
Assume the abc conjecture. Then the equations $x^n+y^n=z^n$ with $n\ge7$ have only finitely many solutions in coprime positive integers, all exponents
together.
\end{proposition}
\begin{proof}
Take $\varepsilon=1$ and $K=K_1$. For a solution, $(x^n,y^n,z^n)$ is an abc triple with $\operatorname{rad}(x^ny^nz^n)=\operatorname{rad}(xyz)\le xyz<z^3$. So
$z^n<K\,(z^3)^2=Kz^6$, that is $z^{n-6}<K$. For $n\ge7$ this gives $z<K$, and then $n<6+\log_2K$ since $z\ge2$. Finitely many $z$ and $n$, hence
finitely many solutions.
\end{proof}

Of course Fermat's last theorem is now a theorem (Wiles, 1995), but the point is the speed: one line from abc. The same argument works for equations
with different exponents.

\begin{lemma}\label{lem:4142}
If $p,q,r$ are positive integers with $\frac1p+\frac1q+\frac1r<1$, then $\frac1p+\frac1q+\frac1r\le\frac{41}{42}$, with equality for $\{p,q,r\}=\{2,3,7\}$.
\end{lemma}
\begin{proof}
Order them as $p\le q\le r$; then $p\ge2$. If $p\ge4$, the sum is at most $\frac34$. If $p=3$, then $q\ge3$; for $q=3$ we need $r\ge4$ and the sum is at most
$\frac1{3}+\frac13+\frac14=\frac{11}{12}$, and for $q\ge4$ it is at most $\frac13+\frac14+\frac14=\frac56$. If $p=2$, then $q\ge3$; for $q=3$ we need
$\frac1r<\frac16$, so $r\ge7$ and the sum is at most $\frac12+\frac13+\frac17=\frac{41}{42}$; for $q=4$, $r\ge5$ and the sum is at most $\frac{19}{20}$; for
$q\ge5$ it is at most $\frac12+\frac15+\frac15=\frac9{10}$. The largest of these bounds is $\frac{41}{42}$.
\end{proof}

\begin{theorem}[abc implies the Fermat--Catalan conjecture]\label{thm:fermatcatalan}
Assume the abc conjecture. Then there are only finitely many triples $(x^p,y^q,z^r)$ of coprime positive integers with $x^p+y^q=z^r$ and
$\frac1p+\frac1q+\frac1r<1$.
\end{theorem}
\begin{proof}
Let $c=z^r$. Then $x<c^{1/p}$, $y<c^{1/q}$ and $z=c^{1/r}$, so $\operatorname{rad}(x^py^qz^r)\le xyz<c^{\frac1p+\frac1q+\frac1r}\le c^{41/42}$ by
Lemma~\ref{lem:4142}. The abc conjecture with $\varepsilon=\frac1{42}$ gives $c<K\,(c^{41/42})^{43/42}=Kc^{1763/1764}$, hence $c<K^{1764}$. There are
finitely many such $c$, and for each finitely many ways to write $c=x^p+y^q$.
\end{proof}

Ten solutions are known:
\[
\begin{gathered}
1^m+2^3=3^2,\quad 2^5+7^2=3^4,\quad 7^3+13^2=2^9,\quad 2^7+17^3=71^2,\quad 3^5+11^4=122^2,\\
17^7+76271^3=21063928^2,\quad 1414^3+2213459^2=65^7,\quad 9262^3+15312283^2=113^7,\\
43^8+96222^3=30042907^2,\quad 33^8+1549034^2=15613^3,
\end{gathered}
\]
the first one for every $m\ge7$ (the notebook checks them all). The conjecture that the list is finite is open; it would follow from abc.

\begin{remark}[more consequences, and what is proved]\label{rem:abcmore}
Elkies showed in 1991 that abc implies Faltings's theorem (a curve of genus at least $2$ has finitely many rational points), with an explicit bound on
the height of the points, which Faltings's proof does not give; Bombieri showed that abc implies Roth's theorem on rational approximations to algebraic
numbers. Granville and Tucker's article in the sources lists many more. Unconditionally, the best result, due to Stewart and Yu (2001), is far weaker:
$c<\exp\bigl(K_\varepsilon\operatorname{rad}(abc)^{1/3+\varepsilon}\bigr)$, an exponential instead of a power.
\end{remark}

\paragraph{Why it is hard.} The conjecture links the prime factorisations of $a$, $b$ and $a+b$, and no known method controls the factorisation of a sum
in terms of the factorisations of its terms: sieves see one number at a time, and linear forms in logarithms (Baker's method, behind Stewart and Yu)
give only exponential bounds. The polynomial proof uses the derivative, which has no good analogue for integers; finding one, an ``arithmetic
derivative'' with the right properties, is part of the motivation of Arakelov theory and of Mochizuki's work.

\begin{remark}[Mochizuki's claimed proof]\label{rem:mochizuki}
In August 2012 Shinichi Mochizuki, of the Research Institute for Mathematical Sciences (RIMS) in Kyoto, posted four long papers on
\emph{inter-universal Teichm\"uller theory}, claiming a proof of abc. Few mathematicians could follow them. In 2018 Peter Scholze and Jakob Stix
spent a week in Kyoto discussing the proof with Mochizuki and Yuichiro Hoshi, and wrote a report explaining why they believe that a key step (Corollary
3.12 of the third paper) does not hold; Mochizuki replied that their objections rest on a misunderstanding. The papers were published in 2021 in
\emph{Publications of RIMS}, a journal edited at Mochizuki's institute. As of 2026, most number theorists do not regard abc as proved, and this book
treats it as open.
\end{remark}

\section{A cousin: the Erd\H os--Straus conjecture}

Egyptian fractions, sums of distinct or repeated unit fractions $\frac1x$, are among the oldest objects of arithmetic. In 1948 Paul Erd\H os and Ernst
Straus asked about the simplest case where a fraction with numerator $4$ must be split.

\begin{center}
\emph{The Erd\H os--Straus conjecture.} For every integer $n\ge2$ there are positive integers $x,y,z$ with $\dfrac4n=\dfrac1x+\dfrac1y+\dfrac1z$.
\end{center}

Like abc, it is a statement about how additive and multiplicative structure meet; unlike abc, it has consequences for nothing, and it is open for the
same reason that it is easy to state: identities settle almost everything, and the remaining cases are the ones where identities cannot reach.

\begin{proposition}\label{prop:esreduce}
If $\frac4m$ is a sum of three unit fractions, so is $\frac4{km}$ for every $k\ge1$. Hence it is enough to prove the conjecture for $n=2$ and for the
odd primes.
\end{proposition}
\begin{proof}
If $\frac4m=\frac1x+\frac1y+\frac1z$, then $\frac4{km}=\frac1{kx}+\frac1{ky}+\frac1{kz}$. Every $n\ge2$ has an odd prime factor or is a power of $2$, hence a
multiple of $2$.
\end{proof}

\begin{lemma}\label{lem:typeI}
Let $a,b,c,d$ be positive integers with $c\mid a+b$, and suppose that $n=\frac{4abcd-a-b}c$ is positive. Then
\[
\frac4n=\frac1{abd}+\frac1{acdn}+\frac1{bcdn}.
\]
\end{lemma}
\begin{proof}
With the common denominator $abcdn$, the right side is $\frac{cn+b+a}{abcdn}=\frac{4abcd}{abcdn}=\frac4n$, because $cn=4abcd-a-b$.
\end{proof}

As $d$ runs over the positive integers, $n$ runs over an arithmetic progression of difference $4ab$. Each choice of $(a,b,c)$ settles a whole residue
class.

\begin{theorem}[Mordell]\label{thm:mordell}
The Erd\H os--Straus conjecture holds for every $n\ge2$ that has a prime factor outside the six residue classes
\[
1,\ 121,\ 169,\ 289,\ 361,\ 529\pmod{840},
\]
and in particular for every $n$ not in these classes. These are the six classes of the squares $1^2,11^2,13^2,17^2,19^2,23^2$ modulo $840$.
\end{theorem}
\begin{proof}
By Proposition~\ref{prop:esreduce} it is enough to settle every prime $p$ outside the six classes, and $p=2$. The prime $2$ is done. If $p\equiv2\pmod3$,
write $p=3k+2$:
\[
\frac4p=\frac1p+\frac1{k+1}+\frac1{(k+1)p},
\]
since $\frac1{k+1}+\frac1{(k+1)p}=\frac{p+1}{(k+1)p}=\frac3p$. All the other cases come from Lemma~\ref{lem:typeI}, with the following $(a,b,c)$ and the
progression of $n$ that they cover:
\[
\begin{array}{lll}
(a,b,c)&n=\frac{4abcd-a-b}c&\text{covers}\\\hline
(1,1,2)&4d-1&n\equiv3\pmod4\\
(1,2,1)&8d-3&n\equiv5\pmod8\\
(1,5,2)&20d-3&n\equiv1\pmod4,\ n\equiv2\pmod5\\
(2,5,1)&40d-7&n\equiv1\pmod8,\ n\equiv3\pmod5\\
(1,21,2)&84d-11&n\equiv1\pmod{12},\ n\equiv3\pmod7\\
(2,21,1)&168d-23&n\equiv1\pmod{24},\ n\equiv5\pmod7\\
(1,14,1)&56d-15&n\equiv1\pmod8,\ n\equiv6\pmod7
\end{array}
\]
(For instance $84d-11\equiv1\pmod{12}$ and $\equiv3\pmod7$, and conversely every $n$ with these two congruences is $\equiv73\pmod{84}$, of the form
$84d-11$ with $d\ge1$.) Now let $p\ge3$ be a prime not covered by the identity mod $3$ and the first two rows. Then $p\equiv1\pmod4$, and not
$\equiv5\pmod8$, so $p\equiv1\pmod8$; $p\equiv1\pmod3$ (as $p\ne3$, which is $\equiv3\pmod4$). If $p\equiv2$ or $3\pmod5$, rows three and four apply. If
$p\equiv3,5,6\pmod7$, the last three rows apply. So $p\equiv1\pmod{24}$, $p\equiv1$ or $4\pmod5$ ($p\ne5$, which is $\equiv5\pmod8$), and $p\equiv1,2$ or
$4\pmod7$ ($p\ne7$, which is $\equiv3\pmod4$). By the Chinese remainder theorem, these are $1\cdot2\cdot3=6$ classes modulo $840$, namely those of
$1^2,11^2,13^2,17^2,19^2,23^2$: each of these squares is $\equiv1\pmod{24}$, and their residues modulo $5$ and $7$ are the six pairs
$(1,1),(1,2),(4,1),(4,2),(1,4),(4,4)$.
\end{proof}

\begin{figure}[ht]
\centering
\fig{ch09-es}
\caption{The $192$ residue classes modulo $840$ that contain primes other than $2,3,5,7$, coloured by the first identity of Theorem~\ref{thm:mordell}
that settles them. Only the six classes of squares, framed in red, are left; there the conjecture is still open.}
\label{fig:es}
\end{figure}

\begin{remark}[the squares, and the evidence]\label{rem:esrest}
The six remaining classes are exactly the squares among the units modulo $840$, and this is no accident: polynomial identities of this kind can never
settle a class of quadratic residues, a fact that goes back to Mordell's book \emph{Diophantine Equations} (1969) and is discussed by Elsholtz and Tao.
Each prime in these classes must be handled by a search, and searches have never failed: the conjecture has been checked for all $n\le10^{17}$ (Salez,
2014). Elsholtz and Tao (2013) showed that the number of solutions for a prime $p$ is, on average over $p$, a power of $\log p$, which makes
counterexamples look very unlikely without saying anything about a given $p$.
\end{remark}

\section{Reading a paper: Granville and Tucker}

Andrew Granville and Thomas Tucker, \emph{It's as easy as abc} (Notices of the AMS 49, 2002), is a short survey written for a general mathematical
audience. Read it in three passes. First its proof of Mason--Stothers, which is the one above, and its explanation of why abc is the natural integer
analogue. Then the list of consequences: each is a few lines of the type of Proposition~\ref{prop:abcfermat}, and you can check the exponents yourself.
Last, the section on Szpiro's conjecture and elliptic curves, which explains where the conjecture came from and is the bridge to the next chapter. The
idea worth stealing is the dictionary between $\mathbb Z$ and $\mathbb C[t]$, which turns hopeless problems about integers into exercises about
polynomials and tells you what to expect.

\section{Open doors}

\begin{itemize}
\item Extend the search of Figure~\ref{fig:abcquality} to $c<10^9$ with a compiled program. How many hits are there, and how does their number grow
with the bound?
\item For each $k\le20$, find the abc triple with $c<10^k$ of largest quality. Does the record ever change after $10^7$?
\item Search for new solutions of $x^p+y^q=z^r$ with $\frac1p+\frac1q+\frac1r<1$ and all terms below $10^{15}$, by meeting in the middle.
\item For the primes $p<10^7$ in the six classes of Theorem~\ref{thm:mordell}, find the smallest $x$ in a solution of $\frac4p=\frac1x+\frac1y+\frac1z$.
How large does it get compared with $p$?
\end{itemize}

\section{Exercises}

\begin{exercise}[$\star$]\label{ex:quality}
Compute the quality of the triples $1+8=9$, $5+27=32$, $1+48=49$ and $1+4374=4375$. Which is a hit?
\end{exercise}

\begin{exercise}[$\star\star$]\label{ex:fermatsquare}
Let $p$ be an odd prime. Show that $p^2$ divides $2^{p(p-1)}-1$, and deduce that $\bigl(1,2^{p(p-1)}-1,2^{p(p-1)}\bigr)$ has
$c/\operatorname{rad}(abc)\ge\frac p2$. (This gives another proof that the $\varepsilon$ is needed.)
\end{exercise}

\begin{exercise}[$\star$]\label{ex:masonsharp}
Show that Theorem~\ref{thm:mason} is sharp: for $a=t^n$, $b=1$, $c=t^n+1$, the two sides are equal.
\end{exercise}

\begin{exercise}[$\star$]\label{ex:charp}
Over the field $\mathbb F_p$, show that $t^p+1=(t+1)^p$, and check that Theorem~\ref{thm:mason} fails for this triple. Which step of the proof breaks?
\end{exercise}

\begin{exercise}[$\star$]\label{ex:pythpoly}
Find coprime nonconstant polynomials with $x^2+y^2=z^2$, and check Theorem~\ref{thm:mason} on them.
\end{exercise}

\begin{exercise}[$\star\star$]\label{ex:explicitabc}
Suppose that every abc triple satisfies $c<\operatorname{rad}(abc)^2$ (an explicit form of abc, still unproved). Deduce Fermat's last theorem for all
exponents $n\ge6$.
\end{exercise}

\begin{exercise}[$\star\star$]\label{ex:catalan}
Assume the abc conjecture. Show that the equation $x^m-y^n=1$ (Catalan's equation) has only finitely many solutions with $x,y\ge2$ and $m,n\ge2$, all
exponents together. (Mih\u ailescu proved in 2002, without abc, that $3^2-2^3=1$ is the only one.)
\end{exercise}

\begin{exercise}[$\star$]\label{ex:essmall}
Write $\frac4n$ as a sum of three unit fractions for $n=2,\dots,10$, and for $n=73$ using Theorem~\ref{thm:mordell}. Is $73$ in one of the six classes?
\end{exercise}

\begin{exercise}[$\star\star$]\label{ex:esfirst}
Find the smallest prime in each of the six classes of Theorem~\ref{thm:mordell}, and a solution for the smallest of them.
\end{exercise}

\begin{exercise}[$\star\star$]\label{ex:esfive}
Show by identities that $\frac5n$ is a sum of three unit fractions whenever $n\equiv0,2,3$ or $4\pmod5$. (Sierpi\'nski conjectured that $\frac5n$ is
always such a sum for $n\ge2$.)
\end{exercise}

\section{Hints}

\begin{hint}{ex:quality}
$\operatorname{rad}(4375\cdot4374)=5\cdot7\cdot2\cdot3$.
\end{hint}
\begin{hint}{ex:fermatsquare}
Euler's theorem modulo $p^2$: $\varphi(p^2)=p(p-1)$.
\end{hint}
\begin{hint}{ex:masonsharp}
The roots of $t^n+1$ are distinct.
\end{hint}
\begin{hint}{ex:charp}
The binomial coefficients $\binom pk$, $0<k<p$, are divisible by $p$. And what is the derivative of $t^p$ in characteristic $p$?
\end{hint}
\begin{hint}{ex:pythpoly}
$(t^2-1)^2+(2t)^2=(t^2+1)^2$.
\end{hint}
\begin{hint}{ex:explicitabc}
Follow Proposition~\ref{prop:abcfermat} with $K=1$.
\end{hint}
\begin{hint}{ex:catalan}
The triple is $(1,y^n,x^m)$, and $\frac1m+\frac1n\le1$; but to beat the radical you need $\frac1m+\frac1n<1$. Treat $m=n=2$ separately.
\end{hint}
\begin{hint}{ex:essmall}
Use the rows of the table; $73=84-11$.
\end{hint}
\begin{hint}{ex:esfirst}
The smallest is $p=1009$. For a solution, try $x=\lceil\frac p4\rceil$ and look for $y$, $z$.
\end{hint}
\begin{hint}{ex:esfive}
Write $n=5k+j$ and subtract $\frac1{k+1}$ from $\frac5n$. For $j=2$, is $(k+1)(5k+2)$ even?
\end{hint}

\section{Solutions}

\begin{solution}{ex:quality}
\idea The radical forgets the exponents.
\writeup $1+8=9$: $\operatorname{rad}=1\cdot2\cdot3=6$, $q=\frac{\log9}{\log6}=1.2263$. $5+27=32$: $\operatorname{rad}=5\cdot3\cdot2=30$,
$q=\frac{\log32}{\log30}=1.0190$. $1+48=49$: $\operatorname{rad}=1\cdot6\cdot7=42$, $q=\frac{\log49}{\log42}=1.0412$. $1+4374=4375$, with $4374=2\cdot3^7$ and
$4375=5^4\cdot7$: $\operatorname{rad}=210$, $q=\frac{\log4375}{\log210}=1.5679$. All four are hits; the last has the second largest quality below
$10^7$.
\end{solution}

\begin{solution}{ex:fermatsquare}
\idea Euler's theorem modulo $p^2$ puts a square in $b$.
\writeup $2$ is prime to $p^2$, and the group $(\mathbb Z/p^2\mathbb Z)^\times$ has order $\varphi(p^2)=p(p-1)$, so $2^{p(p-1)}\equiv1\pmod{p^2}$. Let
$b=2^{p(p-1)}-1$, which is odd. Then $\operatorname{rad}(b)\le\frac bp$ (the prime $p$ appears to the power at least $2$ in $b$ and once in the radical),
and $\operatorname{rad}(1\cdot b\cdot c)=2\operatorname{rad}(b)\le\frac{2b}p<\frac{2c}p$. So $\frac c{\operatorname{rad}(abc)}>\frac p2$, unbounded as $p$ grows.
\end{solution}

\begin{solution}{ex:masonsharp}
\idea All the roots of $c$ are simple, and $a$ contributes one more.
\writeup $abc=t^n(t^n+1)$. The roots are $0$ and the $n$ distinct $n$-th roots of $-1$ (the derivative $nt^{n-1}$ of $t^n+1$ does not vanish at them), so
$n_0(abc)=n+1$ and $\max\deg=n=n_0(abc)-1$.
\end{solution}

\begin{solution}{ex:charp}
\idea In characteristic $p$, a $p$-th power has derivative zero.
\writeup By the binomial theorem, $(t+1)^p=t^p+1+\sum_{0<k<p}\binom pkt^k$, and $p\mid\binom pk$ for $0<k<p$, so the middle terms vanish in $\mathbb F_p$.
With $a=t^p$, $b=1$, $c=(t+1)^p$: $n_0(abc)=2$ (roots $0$ and $-1$) and $\deg c=p>1$. In the proof, $W=ab'-a'b=-pt^{p-1}=0$: the step $W\ne0$ fails,
because $b/a=t^{-p}$ has derivative $0$ without being constant.
\end{solution}

\begin{solution}{ex:pythpoly}
\idea Parametrise Pythagorean triples.
\writeup $x=t^2-1$, $y=2t$, $z=t^2+1$ are coprime ($z-x=2$ and neither has $0$ as a common root with $y$), and
$(t^2-1)^2+4t^2=(t^2+1)^2$. For $a=x^2$, $b=y^2$, $c=z^2$, the roots of $abc$ are $\pm1$, $0$, $\pm i$: $n_0=5$, and $\max\deg=4=n_0-1$. Equality again,
which is why $n=2$ survives in Corollary~\ref{cor:fermatpoly}.
\end{solution}

\begin{solution}{ex:explicitabc}
\idea The proof of Proposition~\ref{prop:abcfermat} with no constant.
\writeup If $x^n+y^n=z^n$ with $x,y,z$ coprime positive integers, then $z^n<\operatorname{rad}(xyz)^2\le(xyz)^2<z^6$, so $n<6$. Hence no solution for
$n\ge6$. (With the classical proofs for $n=3,4,5$, this would have given all of Fermat's last theorem.)
\end{solution}

\begin{solution}{ex:catalan}
\idea The triple $(1,y^n,x^m)$ has a small radical unless $m=n=2$.
\writeup If $m=n=2$, $(x-y)(x+y)=1$ has no solution with $x,y\ge2$. Otherwise $\frac1m+\frac1n\le\frac12+\frac13=\frac56$. Let $c=x^m$; then $y<c^{1/n}$ and
$\operatorname{rad}(1\cdot y^n\cdot x^m)\le xy<c^{1/m+1/n}\le c^{5/6}$. The abc conjecture with $\varepsilon=\frac1{10}$ gives $c<Kc^{(5/6)(11/10)}=Kc^{11/12}$,
so $c<K^{12}$: finitely many $c$, and each $c=x^m$ has finitely many representations with $m\ge2$.
\end{solution}

\begin{solution}{ex:essmall}
\idea Most small $n$ fall into an easy row.
\writeup $\frac42=\frac11+\frac12+\frac12$, $\frac43=\frac11+\frac14+\frac1{12}$, $\frac44=\frac12+\frac13+\frac16$, $\frac45=\frac12+\frac14+\frac1{20}$,
$\frac46=\frac12+\frac17+\frac1{42}$, $\frac47=\frac12+\frac1{15}+\frac1{210}$, $\frac48=\frac13+\frac17+\frac1{42}$, $\frac49=\frac13+\frac1{10}+\frac1{90}$,
$\frac4{10}=\frac13+\frac1{16}+\frac1{240}$. For $73=84\cdot1-11$, the row $(1,21,2)$ with $d=1$ gives
\[
\frac4{73}=\frac1{21}+\frac1{146}+\frac1{3066},
\]
as $\frac{146+21+1}{3066}=\frac{168}{3066}=\frac4{73}$. And $73\equiv73\pmod{840}$ is not one of the six classes.
\end{solution}

\begin{solution}{ex:esfirst}
\idea Walk along each class until a prime appears.
\writeup The smallest primes are $2521$ (class $1$), $1801$ ($121$), $1009$ ($169$), $1129$ ($289$), $1201$ ($361$) and $3049$ ($529$); the numbers
before them in each class are composite (for instance $841=29^2$, $1681=41^2$, $961=31^2$, $289=17^2$, $361=19^2$, $529=23^2$, $1369=37^2$, $2209=47^2$).
For $1009$: with $x=260$, $\frac4{1009}-\frac1{260}=\frac{1040-1009}{1009\cdot260}=\frac{31}{262340}$, and $31=26+5$ with $26$ and $5$ dividing
$262340=2^2\cdot5\cdot13\cdot1009$, so
\[
\frac4{1009}=\frac1{260}+\frac1{10090}+\frac1{52468}.
\]
The trick, splitting the numerator into divisors of the denominator, is how searches work; no identity does it for the whole class.
\end{solution}

\begin{solution}{ex:esfive}
\idea Subtract the largest unit fraction below $\frac5n$ and look at what is left.
\writeup If $n=5k$, $\frac5n=\frac1k=\frac1{3k}+\frac1{3k}+\frac1{3k}$. Otherwise write $n=5k+j$ with $j\in\{2,3,4\}$, so that
$\frac5n-\frac1{k+1}=\frac{5-j}{(k+1)n}$. Let $N=(k+1)n$. For $j=4$: $\frac5n=\frac1{k+1}+\frac1N=\frac1{k+1}+\frac1{2N}+\frac1{2N}$. For $j=3$:
$\frac5n=\frac1{k+1}+\frac1N+\frac1N$. For $j=2$: $N=(k+1)(5k+2)$ is even (if $k$ is even, $5k+2$ is), and $\frac3N=\frac1N+\frac1{N/2}$.
\end{solution}

\section*{Sources and further reading}

\begin{itemize}
\item J.~Oesterl\'e, \emph{Nouvelles approches du ``th\'eor\`eme'' de Fermat}, S\'eminaire Bourbaki 694 (1988), Ast\'erisque 161--162, 165--186; and
D.~W.~Masser, \emph{Open problems}, in \emph{Proceedings of the Symposium on Analytic Number Theory}, Imperial College, London, 1985. For the
history and Goldfeld's phrase: D.~Goldfeld, \emph{Beyond the last theorem}, The Sciences, March/April 1996.
\item A.~Granville and T.~J.~Tucker, \emph{It's as easy as abc}, Notices of the AMS 49 (2002), 1224--1231.
\item W.~W.~Stothers, \emph{Polynomial identities and Hauptmoduln}, Quarterly Journal of Mathematics 32 (1981), 349--370. R.~C.~Mason,
\emph{Diophantine Equations over Function Fields}, Cambridge University Press, 1984. N.~Snyder, \emph{An alternate proof of Mason's theorem}, Elemente
der Mathematik 55 (2000), 93--94.
\item For the records: B.~de Weger, \emph{Algorithms for Diophantine equations}, CWI Tract 65, 1989; the tables of abc triples maintained by Bart de Smit
(Leiden) and the ABC@Home project. The count of $3504$ hits below $10^7$ was computed for this book.
\item N.~D.~Elkies, \emph{ABC implies Mordell}, International Mathematics Research Notices 1991, 99--109. E.~Bombieri, \emph{Roth's theorem and the
abc-conjecture}, preprint, ETH Z\"urich, 1994. C.~L.~Stewart and K.~Yu, \emph{On the abc conjecture II}, Duke Mathematical Journal 108 (2001),
169--181.
\item On the Fermat--Catalan equation and its ten known solutions: F.~Beukers, \emph{The Diophantine equation $Ax^p+By^q=Cz^r$}, Duke Mathematical
Journal 91 (1998), 61--88, and Darmon and Granville, \emph{On the equations $z^m=F(x,y)$ and $Ax^p+By^q=Cz^r$}, Bulletin of the LMS 27 (1995), 513--543.
P.~Mih\u ailescu, \emph{Primary cyclotomic units and a proof of Catalan's conjecture}, Journal f\"ur die reine und angewandte Mathematik 572 (2004),
167--195.
\item S.~Mochizuki, \emph{Inter-universal Teichm\"uller theory I--IV}, Publications of RIMS 57 (2021). P.~Scholze and J.~Stix, \emph{Why abc is still a
conjecture}, report, 2018. D.~Castelvecchi, \emph{Mathematical proof that rocked number theory will be published}, Nature 580 (2020), 177.
\item P.~Erd\H os, a paper in Hungarian on the equation $\frac1{x_1}+\dots+\frac1{x_n}=\frac ab$, Matematikai Lapok 1 (1950), where the
problem of Erd\H os and Straus is stated; Elsholtz and Tao, below, tell its history. L.~J.~Mordell, \emph{Diophantine Equations}, Academic Press, 1969.
C.~Elsholtz and T.~Tao, \emph{Counting the number of solutions to the Erd\H os--Straus equation on unit fractions}, Journal of the Australian
Mathematical Society 94 (2013), 50--105. S.~E.~Salez, \emph{The Erd\H os--Straus conjecture: new modular equations and checking up to
$N=10^{17}$}, arXiv:1406.6307 (2014).
\end{itemize}
